\chapter{Conclusion}\label{chap:conclusion}

\Cref{sec:summary-findings} restates what this thesis found, \cref{sec:answers-rq} answers
the two thesis-level research questions, \cref{sec:conclusion-limitations} states the
limits of those answers, and \cref{sec:future-work} identifies what follows from them. The
evidence is in \cref{chap:results-evaluation}; nothing is re-derived here.


\section{Summary of Findings}\label{sec:summary-findings}

Dataspace connectors are compared on features far more often than on performance, and the
comparisons that do exist rarely hold the identity layer constant
(\cref{sec:gap-analysis}). This thesis measured three \gls{dsp}-conformant connectors on a normalised single-host substrate: the Fraunhofer ISST \gls{edc}, the Factory-X \gls{edc}, and DSP-Native BaSyx. Each was deployed twice, once with its full \gls{dcp} identity stack running and once with identity verification mocked. The difference between
those two deployments is a direct measurement of what identity costs, rather than an
inference from systems that differ in several respects at once.

The campaign comprises 130 measured runs across eight scenarios and six connector--arm
combinations. Five results stand out.

The two \gls{edc} connectors are indistinguishable at the operating point. With identity
enabled a complete transaction takes 4672.5\,ms on the Fraunhofer ISST \gls{edc} and
4649.0\,ms on Factory-X. The 0.5\,\% between them sits well inside their own run-to-run
variation, so it is not read as a difference. DSP-Native BaSyx completes the same
transaction in 2884.5\,ms.

Identity costs a fixed amount per request, and only in the control plane. Measured on the
provider itself it adds 30.0\,ms, 41.3\,ms, and 44.7\,ms. The cost does not grow with the work a request performs. Across a thousandfold increase in catalog size it stays inside a narrow band while the underlying latency rises, so its relative weight falls as requests get larger. The data plane shows no difference at all.

The protocol sets the baseline, not the identity layer. Between 93 and 97\,\% of every
transaction is spent waiting for the two asynchronous state machines to reach their
terminal states. Initiating a negotiation, initiating a transfer, and pulling the data
together account for one to two per cent.

Capacity differs sharply, and the host was never the limit. All six arms track offered
load to 2.5\,tx/s. Beyond that Factory-X delivers about half of what is offered, the
Fraunhofer ISST \gls{edc} plateaus near 4\,tx/s, and BaSyx keeps climbing. Host \gls{cpu}
peaked between 32 and 48\,\%, so each ceiling belongs to the connector and not to the
machine.

The best-looking arm in the campaign is the one that cannot be compared. DSP-Native BaSyx
with identity disabled leads every measurement taken: 814.0\,ms end to end against
4198.0 and 4416.0\,ms for the two \gls{edc} connectors, 39.49\,tx/s at fifty concurrent
clients against 3.34 and 2.24, and the lowest processor, heap, and collector figures of any arm.

That result belongs to the setup rather than to the connector. Its mock validator accepts a token an \gls{edc} consumer cannot produce, so this arm runs without an \gls{edc} consumer at all. Instead, k6 speaks \gls{dsp} to the provider directly and a callback sink absorbs the asynchronous replies (\cref{subsec:mock-arm}). A consumer control plane, its database,
and its status polling are therefore absent from the measurement, and those are most of what the other five arms spend their time on.

The comparison that survives is the provider's own latency, where identity costs 44.7\,ms against 30.0 and 41.3\,ms and all
three connectors sit in the same order of magnitude.

Beyond the measurements, this thesis contributes the instrument that produced them: a
parity-audited benchmark harness, byte-identical across three connector repositories, with
a documented analysis pipeline in which every reported number is traceable to an immutable
result folder (\cref{sec:data-analysis,sec:reproducibility}).


\section{Answers to Research Questions}\label{sec:answers-rq}

\textbf{TRQ1: How do current \gls{dsp}-conformant dataspace connector implementations
compare in end-to-end performance under realistic dataspace data-exchange workloads?}

They are closer than their differing origins suggest at a realistic operating point, and
they diverge under load. At 1\,tx/s the two \gls{edc}-derived connectors are
indistinguishable from each other, and DSP-Native BaSyx completes the same transaction in roughly 60\,\% of their time (\cref{sec:rq1}).

Under rising load the ordering changes. Factory-X is the first to stop converting offered load into completed transactions, at about 2.5\,tx/s. The Fraunhofer ISST \gls{edc} sustains roughly 4\,tx/s before plateauing, and BaSyx continues to deliver at every rung tested (\cref{sec:rq2}). Resource consumption
separates them further, with garbage-collection activity differing by nearly two orders of
magnitude between the \gls{edc} connectors and BaSyx, while none of the three approached
the host's capacity (\cref{sec:rq3}). No connector showed heap growth that would exhaust
its container over a thirty-minute run.

\textbf{TRQ2: Which architectural and implementation-level factors most strongly drive
the observed performance differences?}

Three factors account for most of what was measured, in decreasing order of effect.

The dominant one is the asynchronous protocol design itself. It is common to all
three implementations, so it explains the baseline rather than the differences between
them. Contract negotiation and transfer are state machines whose terminal states a client
must wait for, and that waiting is 93--97\,\% of every transaction (\cref{sec:rq4}). The protocol imposes asynchronous coordination and
therefore shapes the latency budget, but implementation choices determine how quickly the
state transitions are processed; the protocol explains the structure of the delay rather than
a fixed latency value.

The second is where each implementation performs credential verification. On both
\gls{edc} connectors the identity cost concentrates almost entirely in the catalog request, the first protected message of an exchange, while the subsequent asynchronous waits
are barely affected. On DSP-Native BaSyx the same cost appears inside those waits instead.
This pattern suggests the implementations may differ in where or how often credential verification occurs, with the \gls{edc} stack verifying a presentation once per exchange and then proceeding against an established context, and BaSyx validating per message; the current measurements do not establish that mechanism (\cref{sec:x1-overhead}). The per-hop traces collected during the campaign are the evidence that would confirm it directly.

The third is deployment shape. DSP-Native BaSyx collapses control plane and data plane into a single process and serves the transfer synchronously. That is the most plausible account of both its lower baseline latency and its very different garbage-collection profile (\cref{sec:sut,sec:rq3}). That connector's comparison against
the other two is, however, the weakest of the three: its two arms do not share a consumer,
so only its provider-side figures are directly comparable (\cref{subsec:tv-internal}).


\section{Limitations}\label{sec:conclusion-limitations}

Five limits bound what these results support. Everything was measured on one machine, with
the participants communicating over a container network rather than a real one. The
absolute latencies are therefore lower than a distributed deployment would show. The study establishes the ordering of the connectors and the shape of their curves in this single-host setup. Whether those relationships generalise to distributed deployments remains to be validated.

The workload is narrower than a production dataspace: one asset type,
one dataset shape, the negotiate--transfer--pull cycle, and a single membership credential,
so richer usage policies were never exercised. The
thirty-minute endurance window is long enough to show that no component grows fast enough to
exhaust its heap on that timescale, and too short to exclude a slow leak.

No significance test accompanies any comparison. Five repetitions of a run whose internal
transactions are not independent of one another cannot support
one~\citep{hoefler2015scientific}. Differences are therefore judged against the observed
run-to-run range, and reported as inconclusive wherever the ranges of two arms overlap.
\Cref{sec:results-validity} states separately what the reported numbers measure and what
could have biased the comparison between the two arms of a connector.


\section{Future Work}\label{sec:future-work}

The benchmark reported here runs in a controlled and deliberately simplified setting. Four
extensions would move it closer to the conditions a dataspace actually operates in, and each
is a change of experimental design rather than a further run of the present one.

\textbf{Put the participants on separate hosts and separate networks.} A real dataspace is
two organisations reaching each other across a wide-area network, each terminating TLS at its
own boundary and resolving the other's \glspl{did} through infrastructure neither of them
owns. Running the same design with the participants genuinely separated would show which of
the differences reported here survive a real network and which come from co-location. It
would also test the central finding directly: identity costs a fixed amount per request as
long as a \gls{did} document is read from a neighbouring container, and fetching one across a
network may not behave that way. The principles for reproducible measurement in distributed
environments are already established~\citep{papadopoulos2021reproducible}.

\textbf{Grow the number of participants, not only the catalog.} Every measurement here
involved one consumer and one provider. A production dataspace has many of both, and the
component that grows with that number is the catalog: \cref{sec:rq5} varied the assets
advertised by one provider, never the number of providers a consumer must query. Federated
catalog crawling is the operation that scales with participant count, and it already runs as
a separate service in one of the implementations measured here
(\cref{subsec:connector-versions}). Adding participants rather than assets would show whether
connectors that sit close together at two-party scale stay close once discovery is federated.

\textbf{Measure every form of data transfer, not one.} This benchmark moved bytes one way, as
an HTTP pull from a provider endpoint. The \gls{dsp} governs the control plane only, and states that data-plane ``interfaces and interaction patterns are not in scope''~\citep{datapsaceProtocol2025}. How fast a connector moves data is therefore a property of the implementation rather than a guarantee of the protocol. The \gls{edc} ecosystem already carries data-plane implementations beyond
HTTP in separate technology repositories~\citep{eclipseEdcTechnologyAws}, and push transfers
reverse the direction of the movement entirely. Sweeping object storage, message brokers, push
as well as pull, and streaming instead of request--response would show whether identity stays
confined to the control plane when the data plane is not a single synchronous request.

\textbf{Vary the usage policy.} Every negotiation here carried the same minimal policy, and
nothing beyond a membership credential was evaluated. Usage control is what separates a
dataspace connector from an ordinary transfer service, and the policies written in practice
are considerably richer than the one used here~\citep{dam2023policypatterns}. Policy evaluation runs on the critical path of contract negotiation, which is where the identity cost concentrates on the two \gls{edc} connectors (\cref{sec:x1-overhead}). Policy complexity is therefore a plausible second cost of the same shape. Holding the connector and the workload fixed while
varying policy complexity reuses the two-arm design of \cref{sec:gqm-decomposition} directly.
